%
% teil0.tex -- Einleitendes Kapitel: Algebra trifft Geometrie
%
\section{Algebra trifft Geometrie
\label{komplex:section:einleitung}}
\kopfrechts{Algebra trifft Geometrie}
Komplexe Zahlen verbinden zwei Welten, die in der Schulmathematik
oft getrennt unterrichtet werden:
die Algebra mit ihren Symbolmanipulationen
und die elementare Geometrie der Konstruktionen mit Zirkel und Lineal.
\index{Zirkel}%
\index{Linearl}%
\index{Konstruktion mit Zirkel und Lineal}%
Der vorliegende Text zeigt am Beispiel des Satzes von Napoleon,
\index{Napoleon, Satz von}%
wie sich diese Verbindung zu einem schlagkräftigen Werkzeug ausbauen lässt.
Die folgenden Unterabschnitte motivieren das Thema historisch
und geben einen Überblick über den Aufbau der Arbeit.

%
% Motivation und historischer Kontext
%
\subsection{Motivation und historischer Kontext
\label{komplex:subsection:motivation}}
Die komplexen Zahlen sind aus einem rein algebraischen Bedürfnis heraus
entstanden.
Die Lösungsformeln für Gleichungen dritten Grades, die Cardano in der
\index{Cardano, Gerolamo}%
\emph{Ars Magna} von 1545 veröffentlichte,
\index{Ars Magna}%
lieferten in bestimmten Fällen Wurzeln aus negativen Zahlen,
obwohl die Endergebnisse reell waren.
Bombelli erkannte wenige Jahrzehnte später,
dass sich mit diesen scheinbar unsinnigen Ausdrücken konsistent rechnen lässt,
sofern man eine zusätzliche Grösse $i$ mit der Eigenschaft $i^2 = -1$ zulässt.

Die geometrische Deutung der komplexen Zahlen als Punkte einer Ebene
wurde erst zwei Jahrhunderte später von Wessel, Argand und schliesslich
Gauss systematisch ausgearbeitet.
\index{Wessel, Caspar}%
\index{Argand, Jean-Robert}%
\index{Gauss, Carl Friedrich}%
Mit dieser Deutung verbinden komplexe Zahlen zwei Welten:
die Algebra, in der man mit Symbolen rechnet,
und die elementare Geometrie, die mit Zirkel und Lineal Figuren konstruiert.
Eine Drehung in der Ebene wird in der einen Welt durch eine Abbildung
mit Sinus und Kosinus beschrieben,
in der anderen Welt durch eine Multiplikation mit einer einzigen
komplexen Zahl.
Die zweite Beschreibung ist nicht nur kompakter,
sondern lässt sich auch direkt mit Methoden der linearen Algebra
weiterverarbeiten.

Diese doppelte Lesart liefert ein mächtiges Werkzeug für geometrische
Probleme.
Aussagen über Strecken, Winkel und Symmetrien lassen sich in algebraische
Identitäten übersetzen, die man mechanisch nachprüfen kann.
Umgekehrt erhalten algebraische Manipulationen eine anschauliche
Bedeutung, die das Verständnis stützt und Plausibilitätskontrollen
ermöglicht.

%
% Ziel und Aufbau der Arbeit
%
\subsection{Ziel und Aufbau der Arbeit
\label{komplex:subsection:ziel-aufbau}}
Die vorliegende Arbeit zeigt am Beispiel klassischer Konstruktionen,
wie sich elementargeometrische Sätze mit komplexen Zahlen behandeln lassen.
Im Zentrum steht der \emph{Satz von Napoleon}:
Errichtet man über den Seiten eines beliebigen Dreiecks gleichseitige
Dreiecke nach aussen,
dann bilden deren Schwerpunkte selbst wieder ein gleichseitiges Dreieck.
Dieser Satz dient als Leitbeispiel,
an dem alle eingeführten Werkzeuge zusammenspielen:
Polardarstellung, Drehung, Schwerpunktformel und algebraische Identitäten
in $\mathbb{C}$.

Der Aufbau folgt dem natürlichen Lernweg.
Abschnitt~\ref{komplex:section:zahlenebene} fasst die algebraische und
geometrische Darstellung komplexer Zahlen zusammen
und legt damit das Vokabular fest.
Abschnitt~\ref{komplex:section:polar} widmet sich der Polarform und der
Drehung, die für die spätere Konstruktion zentral ist.
Abschnitt~\ref{komplex:section:objekte} zeigt,
wie sich Geraden und Kreise als Niveaumengen einfacher komplexer Ausdrücke
beschreiben lassen.
Abschnitt~\ref{komplex:section:gleichseitig} entwickelt die Konstruktion
gleichseitiger Dreiecke und etabliert damit den Baustein,
aus dem Abschnitt~\ref{komplex:section:napoleon} das Napoleon-Dreieck
zusammensetzt.
In Abschnitt~\ref{komplex:section:beweis} folgt der vollständige
algebraische Beweis des Satzes.
Abschnitt~\ref{komplex:section:rueckblick} fasst zusammen und skizziert
weiterführende Themen.
